\documentclass[11pt]{article}

\usepackage{amsmath}
\usepackage{amssymb}
\usepackage{amsthm}
\usepackage{booktabs}
\usepackage{hyperref}

\hypersetup{
  colorlinks=true,
  linkcolor=black,
  citecolor=black,
  urlcolor=blue,
  pdftitle={Orma: measuring discretion in on-ledger private credit},
  pdfauthor={Alexandre Lemiere, Andrea Gonzalez}
}

\theoremstyle{plain}
\newtheorem{theorem}{Theorem}
\newtheorem{proposition}[theorem]{Proposition}
\newtheorem{lemma}[theorem]{Lemma}
\newtheorem{corollary}[theorem]{Corollary}
\theoremstyle{definition}
\newtheorem{definition}[theorem]{Definition}
\theoremstyle{remark}
\newtheorem*{remark}{Remark}

\DeclareMathOperator{\round}{round}
\DeclareMathOperator{\clamp}{clamp}
\DeclareMathOperator{\idx}{idx}

% Page layout set with plain lengths rather than a geometry package, so the
% preamble stays portable. Values are safe on both letter and A4.
\setlength{\oddsidemargin}{-0.15in}
\setlength{\evensidemargin}{-0.15in}
\setlength{\textwidth}{6.8in}
\setlength{\topmargin}{-0.55in}
\setlength{\headheight}{14pt}
\setlength{\headsep}{16pt}
\setlength{\textheight}{9.4in}
\setlength{\parskip}{0pt}
\setlength{\abovedisplayskip}{6pt plus 2pt minus 2pt}
\setlength{\belowdisplayskip}{6pt plus 2pt minus 2pt}
\setlength{\abovedisplayshortskip}{3pt plus 2pt}
\setlength{\belowdisplayshortskip}{4pt plus 2pt minus 2pt}

% Tighter heading and theorem spacing, done with the kernel's own hooks rather
% than a sectioning package, so the preamble stays to the stated dependencies.
\makeatletter
\renewcommand\section{\@startsection{section}{1}{\z@}%
  {-2.4ex \@plus -1ex \@minus -.2ex}{1.3ex \@plus .2ex}{\normalfont\large\bfseries}}
\renewcommand\subsection{\@startsection{subsection}{2}{\z@}%
  {-2.0ex \@plus -1ex \@minus -.2ex}{0.9ex \@plus .2ex}{\normalfont\normalsize\bfseries}}
\renewcommand\paragraph{\@startsection{paragraph}{4}{\z@}%
  {1.6ex \@plus 1ex \@minus .2ex}{-1em}{\normalfont\normalsize\bfseries}}
\makeatother
\setlength{\topsep}{3pt plus 1pt minus 1pt}
\setlength{\partopsep}{0pt}

\title{Orma: measuring discretion in on-ledger private credit}
\author{Alexandre Lemiere \and Andrea Gonzalez}
\date{13 September 2026}

\begin{document}
\maketitle

\begin{abstract}
\noindent
This is an annex, not a paper. It states the arithmetic behind Orma so a reader can check it
without reading JavaScript. Every formula is implemented in the repository, and every
measured figure was read from live XRP Ledger Devnet state during the De Vinci Blockchain
XRPL Lending Protocol Hackathon, Paris, 12 and 13 September 2026, against \texttt{rippled}
3.4.0-rc5. The central claim is combinatorial: total first-loss capital consumed across a
fixed set of defaults depends on the order in which they are declared, that order is chosen
unilaterally by the party whose capital is consumed, and a closed form makes the size of the
choice computable from public state. On the measured facility it is worth $0.20$ XRP out of
$0.70$.
\end{abstract}

\section{The setting}

\subsection{The objects}

A \emph{Vault} (XLS-65) holds pooled assets and issues a share token. To carry lending it
must be closed-ended: \texttt{VaultKind} set to $1$, with a \texttt{SubscriptionDate} and a
\texttt{RedemptionDate}. Without those three fields \texttt{LoanBrokerSet} returns
\texttt{tecNO\_PERMISSION}, and \texttt{VaultKind} is immutable, so this is decided at
creation. The vault runs through Subscription, Investment and Redemption, and
\texttt{VaultWithdraw} returns \texttt{tecTOO\_SOON} for the whole of the Investment phase
regardless of available liquidity: inability to withdraw then is a protocol fact, not a
distress signal, and nothing below treats it as one.

A \emph{LoanBroker} (XLS-66) attaches to exactly one vault. It draws capital from the vault,
originates \emph{Loan} objects, and posts \emph{first-loss capital}, called cover, which is
liquidated ahead of the vault's own assets when a loan is written off. The broker declares
the credit events: impairment, un-impairment, default. The share token is an
\texttt{MPTokenIssuance} whose \texttt{OutstandingAmount} is the unit count and whose
metadata the issuer may write, which is what lets a share point at its own valuation
(Section~6).

\subsection{The identity that generates the problem}

Two protocol facts combine into one asymmetry. First, only the \texttt{LoanBroker} owner may
impair, un-impair or default a loan. Second, the \texttt{LoanBroker} owner is necessarily
the vault owner: no other account can create the broker, and XLS-75 permission delegation
returns \texttt{temMALFORMED} across the entire lending suite, so the authority cannot be
split off to a third party even voluntarily.

Therefore the party who decides \emph{when} a loss is recognised, the party who decides
\emph{in what order} a set of losses is realised, and the party whose own first-loss capital
is consumed when they are, are one account. Write $\tau$ for the timing choice and $\sigma$
for the ordering choice. Section~2 measures what $\tau$ is worth, Section~4 what $\sigma$ is
worth. Both move value from unit holders to the manager, and neither is visible to a
consumer using the standard tooling. Orma measures both from public state, publishes the
result as a native XLS-47 Oracle object and as a pointer inside the share token, and lets an
independent third party gate subscriptions on it through XLS-70 credentials. The enforcement
layer contains no arithmetic and is not analysed here.

\subsection{Notation}

\begin{center}
\begin{tabular}{clc}
\toprule
symbol & meaning & ledger field \\
\midrule
$A$ & total assets held by the vault & \texttt{Vault.AssetsTotal} \\
$A_v$ & assets not currently lent out & \texttt{Vault.AssetsAvailable} \\
$L$ & loss recognised but not yet deducted & \texttt{Vault.LossUnrealized} \\
$S$ & units outstanding & \texttt{MPTokenIssuance.OutstandingAmount} \\
$D$ & broker's total drawn debt & \texttt{LoanBroker.DebtTotal} \\
$C$ & first-loss capital posted & \texttt{LoanBroker.CoverAvailable} \\
$c_{\min}$ & minimum cover rate & \texttt{LoanBroker.CoverRateMinimum} \\
$c_{\mathrm{liq}}$ & liquidation cover rate & \texttt{LoanBroker.CoverRateLiquidation} \\
$p_i$ & principal of exposure $i$ & \texttt{Loan.PrincipalOutstanding} \\
\bottomrule
\end{tabular}
\end{center}

\paragraph{Units.} Monetary quantities are integer drops, $1$ XRP $= 10^{6}$ drops, carried
on the wire as decimal strings of up to $19$ significant digits. All arithmetic runs at $40$
significant digits with \texttt{decimal.js}; no monetary value passes through
\texttt{Number()} or \texttt{parseFloat}, because a $53$-bit mantissa cannot represent every
value the protocol admits. $c_{\min}$ and $c_{\mathrm{liq}}$ are in units of $10^{-5}$, so a
stated \texttt{10000} means $0.10$: not $10000$, and not $100$.

\paragraph{Accounting basis.} All facilities here run \texttt{LEVersion} $=1$, cash basis.
The consequence that matters is that origination moves \texttt{DebtTotal}, not
\texttt{AssetsTotal}: lending does not reduce the assets the vault reports holding. This is
load-bearing for Proposition~\ref{prop:empty}.

\subsection{The measured facilities}

Three live Devnet facilities are used below. Their ids are collected here so the text can
name them in words and every figure quoted later can be re-derived from public state.

\begin{center}
\footnotesize
\begin{tabular}{ll}
\toprule
object & id \\
\midrule
Meridian Trade Finance I, vault & \texttt{864C5A2DDCED78C6198B0703169D533B747A56E5F9398D832033F79E37E8C835} \\
Meridian, impairing transaction & \texttt{075FE6D2E0F29919AF477A2A8F581A680805A006138BB967D8434611E49229C3} \\
Kestrel Bridge Financing II, vault & \texttt{24EAA01AD4CE70D8ABB4ACB4255DC4C216B8AF3B9E7315DB8052BFC1E5810089} \\
Kestrel, loan broker & \texttt{8DB104980C31CBE9F196508E7D748E8C31FB783CD4E0A3620AE0B476CBDEA3FA} \\
Kestrel, first default & \texttt{FB23EC5313F85FA157D56966E0B0B1665068905988D925C70BB0E7C35C3B221B} \\
Kestrel, second default & \texttt{DE49DF0D6136B96AD28B253CD32100E5715078FD2D0AE6D1DB747AFECB803D1F} \\
Calder Structured Credit III, broker & \texttt{1C160EFFCA8208DFD8CE9B065A3C3C41CCA80758FBC3099BE27AACFBD1478173} \\
Calder, share token issuance & \texttt{000000014D667775372D5B78E07FFF294678C7F9CE82AFBC} \\
Kestrel, broker and vault owner & \texttt{rD26gw5goDqU784caePMoA82givcJBjStk} \\
\bottomrule
\end{tabular}
\end{center}

\section{The two readings of net asset value}

\subsection{Definitions}

\begin{definition}
The \emph{naive} unit value, which is what a tool reading the headline asset figure
computes, and the \emph{held} unit value, which is what a holder can redeem against, are
\[
  \mathrm{NAV}_{\mathrm{naive}} = \frac{A}{S},
  \qquad
  \mathrm{NAV}_{\mathrm{held}} = \frac{A - L}{S},
\]
both defined as $0$ when $S = 0$.
\end{definition}

They agree exactly until a loss is recognised and diverge the instant one is. Divergence is
reported in basis points,
\[
  \delta \;=\; \round\!\left( 10^{4} \cdot
  \frac{\mathrm{NAV}_{\mathrm{naive}} - \mathrm{NAV}_{\mathrm{held}}}{\mathrm{NAV}_{\mathrm{naive}}} \right)
  \;=\; \round\!\left( 10^{4} \cdot \frac{L}{A} \right),
\]
to the nearest integer basis point, halves away from zero. The second equality holds for
$S \neq 0$, $A \neq 0$, and shows that divergence does not depend on the unit count at all:
it is the recognised loss as a fraction of the reported book.

\begin{remark}
$\mathrm{NAV}_{\mathrm{naive}}$ can \emph{rise} while a vault is emptied, because whenever
$L > 0$ a withdrawal burns units faster than it removes assets.
\end{remark}

\subsection{Why the naive reading is not merely lazy}

The obvious objection is that any competent indexer nets the recognised loss, so the naive
reading is a straw man. It is not, and the reason is a property of the ledger's metadata
rather than of anyone's code.

\begin{proposition}[The empty change set]
\label{prop:empty}
Let $V$ be a Vault under \texttt{LEVersion} $=1$ whose \texttt{LossUnrealized} is $0$ just
before a \texttt{LoanManage} transaction with the impair flag set, and let that transaction
succeed and set \texttt{LossUnrealized} to some $L' > 0$. Then the \texttt{PreviousFields}
object of $V$'s node in that transaction's metadata is empty.
\end{proposition}

\begin{proof}
The metadata for a modified node lists in \texttt{PreviousFields} exactly those fields whose
value before the transaction differs from their value after it, and omits any field whose
previous value was the default for that field's type. For the amount-valued fields of a
Vault that default is $0$.

Take each field in turn. \texttt{AssetsTotal} is unchanged: under cash-basis accounting
origination recorded the principal in \texttt{DebtTotal}, and recognising a loss on an
outstanding loan writes no asset movement. \texttt{AssetsAvailable} is unchanged: nothing is
lent, repaid, deposited or withdrawn. The share issuance is untouched, so $S$ is unchanged,
and every remaining field of the Vault is static configuration. The single field whose value
differs is \texttt{LossUnrealized}, moving from $0$ to $L'$; but its previous value is $0$,
the type default, so by the omission rule it is not emitted. No other field differs, hence
\texttt{PreviousFields} is empty.
\end{proof}

\begin{corollary}
\label{cor:blind}
Let $f$ be any function of the \texttt{PreviousFields} object alone. Then $f$ takes the same
value on the metadata of the impairing transaction as on the metadata of a transaction that
changed nothing about the vault. An indexer maintaining vault state by applying metadata
diffs therefore continues to report
$\mathrm{NAV}_{\mathrm{naive}} = \mathrm{NAV}_{\mathrm{held}}$ indefinitely after the first
impairment, and reports no event at all.
\end{corollary}

The scope is narrow and still damaging. It applies to the \emph{first} impairment of a
previously healthy vault; on the second, $L$ is already non-zero, the field is emitted
normally, and a diffing consumer catches up. The blind spot is exactly the transition from
healthy to impaired, which is the one a lender cares about most. The remedy costs one round
trip: Orma never diffs \texttt{PreviousFields} but polls each vault object every $4$
seconds, re-reads it outright and recomputes both readings from the current values.
\texttt{FinalFields} in the same metadata also carries the correct post-state, so the
failure is specific to the diff.

\subsection{Measured}

Facility \emph{Meridian Trade Finance I}, impairing transaction as listed in Section~1.4,
result \texttt{tesSUCCESS}.

\begin{center}
\begin{tabular}{lrr}
\toprule
field & healthy, before & final, after \\
\midrule
\texttt{AssetsTotal} & $51{,}000{,}000$ & $51{,}000{,}000$ \\
\texttt{AssetsAvailable} & $41{,}000{,}000$ & $41{,}000{,}000$ \\
\texttt{LossUnrealized} & $0$ & $10{,}000{,}000$ \\
\midrule
\multicolumn{3}{l}{\texttt{PreviousFields} emitted for the Vault node: \texttt{\symbol{123}\symbol{125}}} \\
\bottomrule
\end{tabular}
\end{center}

With $S = 51{,}000{,}000$ units,
\[
  \mathrm{NAV}_{\mathrm{naive}} = \frac{51{,}000{,}000}{51{,}000{,}000} = 1.000000,
  \qquad
  \mathrm{NAV}_{\mathrm{held}} = \frac{41{,}000{,}000}{51{,}000{,}000} = 0.803922,
\]
the second rounded to nearest at the $6$ decimal places the valuation is reported to. The divergence
is $10^4 \cdot 10{,}000{,}000 / 51{,}000{,}000 = 1960.7843\ldots$, so $\delta = 1961$ bp: the
reported unit value overstates recoverable value by $19.61\%$. Two readers pointed at the
same transaction on the same ledger at the same instant return $1.000000$ and $0.803922$.

\section{Cover liquidated on a default}

When the broker declares a default on a loan of principal $p$, the ledger liquidates
\begin{equation}
\label{eq:cover}
  T(p) \;=\; \min\Big( \big\lceil D \cdot c_{\min} \cdot c_{\mathrm{liq}} \big\rceil,\; p,\; C \Big)
\end{equation}
of posted cover. Three properties of \eqref{eq:cover} matter and none is obvious from the
prose specification.

\paragraph{The double product.} Both rates are $10^{-5}$ units, so both are divided by
$10^5$. At $c_{\min} = c_{\mathrm{liq}} = 10000$ the coefficient is not $0.10$ but
$0.10 \times 0.10 = 0.01$. Reading one of the two as a plain multiplier overstates
liquidatable cover tenfold.

\paragraph{The ceiling.} The product is rounded up to whole drops, not truncated. At the
scales used here the argument is an exact integer and the ceiling is inert, but it is
implemented because the ledger implements it and it becomes visible at odd rates.

\paragraph{The base.} The base is the broker's \emph{total} outstanding book $D$, not the
principal $p$ of the exposure that defaulted: a default is charged against the size of the
whole book at the instant it is declared. That choice is what makes Section~4 non-trivial,
and it also makes concentration and under-collateralisation one defect seen from two sides,
since a book of one very large loan liquidates no more cover on its default than a book of
many small ones liquidates on the first of them.

\paragraph{Worked, from live state.} The \emph{Calder Structured Credit III} broker, with
$D = 10{,}000{,}000$ drops, $C = 5{,}000{,}000$ drops and
$c_{\min} = c_{\mathrm{liq}} = 10000$:
\[
  c = 0.10 \times 0.10 = 0.01,
  \qquad
  \lceil 10{,}000{,}000 \times 0.01 \rceil = 100{,}000 \text{ drops} = 0.10 \text{ XRP}.
\]
The broker posted $5$ XRP of cover, of which $0.10$ XRP is reachable by any single default
at the current book size: the stranded fraction
$\max(0, C - \lceil D c_{\min} c_{\mathrm{liq}} \rceil)/C$ is $0.98$. The separately defined
cover \emph{requirement} $\lceil D c_{\min} \rceil = 1{,}000{,}000$ drops is met with room to
spare, so a consumer checking the requirement alone concludes the facility is well covered.

\section{The ordering result}

\subsection{Statement}

Fix a broker and write $c = c_{\min} \cdot c_{\mathrm{liq}}$. Let $D_0$ be the book at the
first default, $p_1, \ldots, p_k$ the principals of $k$ exposures that will all be
defaulted, and $\sigma$ a permutation of $\{1, \ldots, k\}$ giving the order in which the
broker declares them. \texttt{DebtTotal} is decremented by $p$ as each default settles, so
the book facing step $i$ is $D_{i-1} = D_0 - \sum_{j < i} p_{\sigma(j)}$. Write
$T(\sigma, k)$ for total cover consumed across the $k$ declarations, under three regularity
assumptions that hold on the measured facility and that Section~7 revisits.

\begin{itemize}
\itemsep2pt
\item[(A1)] $c \, D_{i-1}$ is an integer number of drops at every step, so the ceiling in
\eqref{eq:cover} is inert.
\item[(A2)] Cover never runs out: $C \geq T(\sigma, k)$ for every $\sigma$ considered.
\item[(A3)] Each principal is at least the step demand, $p_{\sigma(i)} \geq c \, D_{i-1}$,
so the per-step $\min$ never binds on $p$.
\end{itemize}

\begin{lemma}[Closed form]
\label{lem:closed}
Under (A1) to (A3),
$\displaystyle T(\sigma, k) = c \left[\, k D_0 - \sum_{i=1}^{k} (k - i)\, p_{\sigma(i)} \right].$
\end{lemma}

\begin{proof}
Under the assumptions the amount liquidated at step $i$ is exactly $c \, D_{i-1}$, so
\[
  T(\sigma, k) = \sum_{i=1}^{k} c \, D_{i-1}
  = c \sum_{i=1}^{k} \Big( D_0 - \sum_{j<i} p_{\sigma(j)} \Big)
  = c \left[ k D_0 - \sum_{i=1}^{k} \sum_{j<i} p_{\sigma(j)} \right].
\]
Exchange the order of summation. The term $p_{\sigma(j)}$ appears once for each $i$ with
$j < i \leq k$, that is $k - j$ times, so
$\sum_{i=1}^{k} \sum_{j<i} p_{\sigma(j)} = \sum_{j=1}^{k} (k-j)\, p_{\sigma(j)}$, which is
the claim after renaming the index.
\end{proof}

Only the second term depends on $\sigma$, and it enters with a minus sign, so minimising
total cover consumed is the same problem as maximising
$\Sigma(\sigma) = \sum_i (k-i)\, p_{\sigma(i)}$: a sum of the principals against the fixed,
strictly decreasing weight vector $(k-1, k-2, \ldots, 1, 0)$.

\begin{theorem}[Ordering]
\label{thm:order}
Under (A1) to (A3), $T(\sigma, k)$ is minimised when $\sigma$ declares the exposures in
non-increasing order of principal, largest first, and maximised when it declares them in
non-decreasing order, smallest first. The extremal values are
\[
  T_{\min} = c \left[ k D_0 - \sum_{i=1}^{k} (k-i)\, p_{(i)}^{\downarrow} \right],
  \qquad
  T_{\max} = c \left[ k D_0 - \sum_{i=1}^{k} (k-i)\, p_{(i)}^{\uparrow} \right],
\]
where $p^{\downarrow}$ and $p^{\uparrow}$ are the principals sorted decreasing and
increasing.
\end{theorem}

\begin{proof}
By Lemma~\ref{lem:closed} it suffices to show $\Sigma$ is maximised by the decreasing
arrangement and minimised by the increasing one. This is the rearrangement inequality; the
exchange argument is given in full because it also yields uniqueness up to ties.

Suppose $\sigma$ is not non-increasing in principal. Then there are adjacent positions $i$
and $i+1$ with $a = p_{\sigma(i)} < p_{\sigma(i+1)} = b$. Let $\sigma'$ be $\sigma$ with
those two transposed. Every other position contributes identically, so
\[
  \Sigma(\sigma') - \Sigma(\sigma)
  = \big[(k-i)\,b + (k-i-1)\,a\big] - \big[(k-i)\,a + (k-i-1)\,b\big]
  = b - a \;>\; 0 ,
\]
because the two weights differ by exactly $1$. Every adjacent transposition that repairs an
inversion therefore strictly increases $\Sigma$ and strictly decreases $T$. A permutation of
$k$ elements has finitely many inversions and each such transposition removes one, so the
process terminates, and it can terminate only at an arrangement with no inversions, that is
the non-increasing one. Hence the decreasing order attains the maximum of $\Sigma$ and no
other order class does. The minimising case is identical with the inequality reversed.
\end{proof}

\begin{corollary}[The spread]
\label{cor:spread}
$T_{\max} - T_{\min} = c \sum_{i=1}^{k} (k-i) \big( p_{(i)}^{\downarrow} - p_{(i)}^{\uparrow} \big)$.
For $k = 2$ this collapses to $c\,( p_{\max} - p_{\min} )$.
\end{corollary}

\begin{corollary}[Stated plainly]
\label{cor:plain}
The party who chooses $\sigma$ is the \texttt{LoanBroker} owner, who is necessarily the vault
owner. The capital consumed by $T(\sigma, k)$ is that same party's posted cover, and every
drop of cover not consumed is a drop of loss falling instead on the vault's unit holders.
The party choosing the order is exactly the party whose capital the choice spares, and the
amount spared, against the investor-optimal order, is $T_{\max} - T(\sigma, k)$.
\end{corollary}

\subsection{Measured on Devnet}

Facility \emph{Kestrel Bridge Financing II}, ids in Section~1.4. Two loans of $30$ XRP and
$10$ XRP, $D_0 = 40$ XRP, $C = 10$ XRP posted, $c_{\min} = c_{\mathrm{liq}} = 10000$ so
$c = 0.01$. The broker declared the large loan first. Both orders, in XRP:

\begin{center}
\begin{tabular}{llrrrr}
\toprule
order & step & book $D_{i-1}$ & principal $p_{\sigma(i)}$ & $\lceil c D_{i-1} \rceil$ & cumulative \\
\midrule
largest first    & 1 & $40$ & $30$ & $0.40$ & $0.40$ \\
(observed)       & 2 & $10$ & $10$ & $0.10$ & $\mathbf{0.50}$ \\
\addlinespace
smallest first   & 1 & $40$ & $10$ & $0.40$ & $0.40$ \\
(counterfactual) & 2 & $30$ & $30$ & $0.30$ & $\mathbf{0.70}$ \\
\bottomrule
\end{tabular}
\end{center}

The closed form agrees. With $k = 2$, $D_0 = 40$, $c = 0.01$ and weights $(1,0)$,
\[
  T(\text{largest first}) = 0.01\,[\,80 - 30\,] = 0.50,
  \qquad
  T(\text{smallest first}) = 0.01\,[\,80 - 10\,] = 0.70,
\]
and Corollary~\ref{cor:spread} gives the spread directly as $0.01 \times (30 - 10) = 0.20$
XRP without computing either endpoint. The ledger agrees too: reconstructed from public
transaction history alone, across the two default transactions of Section~1.4
\texttt{CoverAvailable} moved
$10{,}000{,}000 \rightarrow 9{,}600{,}000 \rightarrow 9{,}500{,}000$ drops, consuming
$400{,}000$ and $100{,}000$ drops, total $500{,}000$ drops $= 0.50$ XRP. Three independent
computations agree: the step-by-step simulation, the closed form, and the ledger itself.

Identical losses were realised in both rows; the only difference is the order of two
transactions eight seconds apart. The $0.20$ XRP was borne by the vault's investors, and
the party who chose the order is the party it spared. Measured against the $0.50$ XRP the
first-loss capital did absorb, the investor-optimal order would have absorbed $40\%$ more;
measured against that $0.70$ XRP, the choice withheld $28.6\%$ of it.

\subsection{Sequence score}

Observed consumption is placed between the two extremes:
\[
  \phi \;=\; \frac{T_{\mathrm{observed}} - T_{\min}}{T_{\max} - T_{\min}} \;\in\; [0, 1],
\]
defined as $1$ when $T_{\max} = T_{\min}$, which happens when all principals are equal or
fewer than two defaults have occurred. $\phi = 0$ is the investor-worst sequence and
$\phi = 1$ the investor-best. For Kestrel,
$\phi = (0.50 - 0.50)/(0.70 - 0.50) = 0.0000$. A separate ordinal conduct score starts at
$100$ and penalises $-35$ for $\phi < 0.25$, $-15$ for $0.25 \leq \phi < 0.75$, $-20$ per
loss written off with no prior impairment and $-10$ per withdrawal of cover; Kestrel scores
$25$, grade E.

Reconstructing this history needs a join that is easy to miss. A \texttt{LoanManage}
transaction carries no \texttt{LoanBrokerID}, and an impairment does not touch the
\texttt{LoanBroker} object at all, so filtering \texttt{account\_tx} on the broker object
silently drops every impairment and returns a history in which no loss was ever signalled
before it was taken. The \texttt{Loan} node in the same metadata carries
\texttt{FinalFields.LoanBrokerID}, which closes the join at no extra request. Before that
fix, a rule written to reward early disclosure was penalising it.

\section{The scoring model}

\subsection{Five factors, no weights}

Five quantities are measured from ledger state, each mapped independently onto a $20$-step
ordinal ladder. No weighted average is taken at any stage: a weighted average lets a strong
factor pay for a broken one, which is precisely the failure mode a credit assessment must
not have.

\begin{description}
\itemsep2pt
\item[Liquidity.] $A_v / A$.
\item[First-loss adequacy.] Liquidatable cover against the exposure it must absorb,
\[
  \mathrm{adequacy} =
  \frac{\min\big( \lceil D c_{\min} c_{\mathrm{liq}} \rceil,\; C \big)}
       {\max\big( L,\; P_{\mathrm{distressed}} \big)},
\]
defined as $1$ when the denominator is zero. The denominator takes the \emph{larger} of the
recognised loss and the principal of exposures that are overdue, impaired or defaultable.
Dividing by $L$ alone would grade a vault concealing a large overdue exposure as perfectly
covered. Concealment must never improve a score; the $\max$ is what enforces that.
\item[Concentration.] $\max_i (p_i) / D$, banded on $1$ minus that value so higher is better.
\item[Recognition lag.] Seconds an exposure has been past due with nothing declared. This is
the direct measurement of the timing lever $\tau$.
\item[Redemption cliff.] Projected shortfall of claims over liquidity at redemption,
\[
  \mathrm{shortfall} = \max\Big( 0,\; (A - L) - \big( A_v + R_{\mathrm{perf}}
  + \mathbf{1}[P_{\mathrm{distressed}} > 0] \cdot C_{\mathrm{rec}} \big) \Big),
\]
with $R_{\mathrm{perf}}$ the total contractual repayment of performing loans and
$C_{\mathrm{rec}}$ the recoverable cover
$\min(\lceil D c_{\min} c_{\mathrm{liq}} \rceil, C)$. \texttt{rippled} refuses a
\texttt{LoanSet} whose term overruns the vault's \texttt{RedemptionDate}, so every
performing loan is contractually repaid before redemption and a projected shortfall is
exactly the non-performing book net of recoverable cover. Being fully lent is therefore not
a distress signal here.
\end{description}

The bands are stated in full so they can be disputed. Each threshold is a floor: the first
one the value meets or exceeds gives the grade, and a value below all of them grades
\texttt{D}. Recognition lag is banded on seconds rather than a ratio: $0$ is AAA, under $60$
is A, under $300$ is BBB, under $900$ is BB, beyond that CCC.

\begin{center}
\small
\begin{tabular}{lll}
\toprule
factor & quantity banded & thresholds and grades \\
\midrule
Liquidity & $A_v/A$ &
  $0.95$ AAA, $0.85$ A, $0.70$ BBB, $0.50$ BB, $0.25$ B, $0.05$ CCC \\
First-loss adequacy & adequacy &
  $1$ AA, $0.50$ A-, $0.20$ BBB-, $0.05$ BB-, $0.01$ B- \\
Concentration & $1 - \max_i(p_i)/D$ &
  $0.80$ AA, $0.60$ A-, $0.40$ BBB-, $0.20$ BB- \\
Redemption cliff & $1 - \mathrm{shortfall}\%/100$ &
  $1$ AAA, $0.98$ AA, $0.90$ A, $0.75$ BBB, $0.50$ BB, $0.25$ B, $0.10$ CCC \\
\bottomrule
\end{tabular}
\end{center}

\subsection{Composite by ordinal notching}

Let $\mathcal{L} = [\text{AAA}, \text{AA+}, \ldots, \text{D}]$ with $|\mathcal{L}| = 20$,
ordered best to worst, and let $\idx(g)$ be the zero-based position of grade $g$. Define
\[
  \mathrm{notch}(g, \Delta) = \mathcal{L}\big[\, \clamp(\idx(g) - \Delta,\; 0,\; 19) \,\big],
  \qquad
  \mathrm{numeric}(g) = \round\!\Big( 100 \big( 1 - \idx(g)/19 \big) \Big),
\]
so a negative $\Delta$ moves the grade down the ladder and $\mathrm{numeric}$ projects it to
a sortable number. The composite \emph{anchors} on the redemption cliff, because for a
fixed-term facility the headline question is whether claims can be met at redemption, then
notches from there. The realised-destruction rows are mutually exclusive; the rest
accumulate. Every notch is emitted with its rule, so a grade is reproducible from the
published trace rather than taken on trust.

\begin{center}
\small
\begin{tabular}{lrlr}
\toprule
realised destruction $\rho$ & $\Delta$ & other conditions & $\Delta$ \\
\midrule
$\rho \geq 75\%$ & $-14$ & recognition lag $> 300$ s & $-2$ \\
$\rho \geq 50\%$ & $-12$ & recognition lag $> 0$ s and $\leq 300$ s & $-1$ \\
$\rho \geq 25\%$ & $-8$ & largest exposure $\geq 75\%$ of the book & $-1$ \\
$\rho \geq 10\%$ & $-5$ & liquidatable cover $< 5\%$ of the exposure & $-1$ \\
$\rho \geq 2\%$ & $-3$ & liquidity band B or worse, in Redemption & $-2$ \\
$\rho > 0\%$ & $-1$ & & \\
\bottomrule
\end{tabular}
\end{center}

\subsection{Realised capital destruction, and the denominator}

Each of the five factors measures current \emph{exposure}, and a realised loss leaves none
behind: the write-off removes the asset, the provision is released, and the book reads
clean. Without a memory term a facility that defaulted on four fifths of its loans scores
identically to one that never lost anything, the instant the write-off settles. That was
observed on Devnet, not hypothesised. The term is
\[
  \rho \;=\; \max\left( 0,\; \frac{S - A}{S} \right),
\]
measured against $A$ and explicitly \emph{not} against $A - L$.

The quantity subtracted is $A$, and that choice carries the whole substance of the rule. An
unrealised loss $L$ is a provision against an asset the vault still holds and may still
recover: it has not destroyed capital, it has disclosed that capital may be destroyed.
Counting $L$ in $\rho$ would notch a manager down at the moment they disclose, punishing the
exact behaviour the recognition-lag factor exists to reward and rewarding a manager who says
nothing. Only a write-off, which removes the asset from \texttt{AssetsTotal} outright, is
permanent destruction, and only that is counted. Par is taken as $1.0$, so $S$ is capital
subscribed measured in asset units; Section~7 states when that holds.

\subsection{Two worked facilities}

\paragraph{Calder Structured Credit III, composite AA.} Live state: $A = 51{,}000{,}000$,
$A_v = 41{,}000{,}000$, $L = 10{,}000{,}000$, $S = 51{,}000{,}000$, $D = 10{,}000{,}000$,
$C = 5{,}000{,}000$, one impaired loan of principal $10{,}000{,}000$. Claims at redemption
are $A - L = 41{,}000{,}000$ and liquidity at redemption is
$41{,}000{,}000 + 0 + 100{,}000 = 41{,}100{,}000$, so the shortfall is $0$ and the anchor is
AAA. Then $\rho = 0$ because $A \geq S$, no notch. Recognition lag is $0$ because the
exposure is already declared impaired, no notch. Concentration is
$10^{7}/10^{7} = 1.00 \geq 0.75$, one notch. Adequacy is
$\min(10^{5}, 5 \cdot 10^{6}) / \max(10^{7}, 10^{7}) = 0.0100 < 0.05$, one notch. So
$\text{AAA} \rightarrow \text{AA+} \rightarrow \text{AA}$, and
$\mathrm{numeric}(\text{AA}) = \round(100(1 - 2/19)) = 89$.

\paragraph{Kestrel Bridge Financing II, composite B.} The same broker as Section~4.2, after
both defaults settled: $A = 10{,}500{,}000$ drops against $S = 50{,}000{,}000$ units
subscribed. The shortfall is $0$, so the anchor is again AAA, and on the five factors alone
the facility grades AAA: there is no exposure left to measure, because it was destroyed
rather than provisioned.
\[
  \rho = \frac{50{,}000{,}000 - 10{,}500{,}000}{50{,}000{,}000} = 0.79 = 79\% \geq 75\%
  \;\Longrightarrow\; \Delta = -14 .
\]
$\idx(\text{AAA}) = 0$, so the composite is $\mathcal{L}[14] = \text{B}$ with
$\mathrm{numeric}(\text{B}) = \round(100(1 - 14/19)) = 26$. Fourteen notches on the strength
of one term. Without it, the most damaged facility in the book would sort first.

\section{Valuing a pledge}

A holder may pledge units to a second lender, who then holds a token and nothing else. For
$u$ units,
\[
  V_{\mathrm{held}} = u \cdot \mathrm{NAV}_{\mathrm{held}},
  \qquad
  V_{\mathrm{reported}} = u \cdot \mathrm{NAV}_{\mathrm{naive}},
  \qquad
  \Omega = V_{\mathrm{reported}} - V_{\mathrm{held}} ,
\]
with $\Omega$ the overstatement the second lender carries if they price the collateral off
the reported figure. By the identity in Section~2,
$\Omega = u L / S = (\delta/10^4) \cdot V_{\mathrm{reported}}$ up to the rounding of
$\delta$. The lender reaches $\mathrm{NAV}_{\mathrm{held}}$ from the token alone in five
steps, none requiring trust in whoever computed it: read the \texttt{MPTokenIssuance},
decode its metadata, check conformance, find the valuation pointer the issuer declared,
follow it. The instrument carries the address of its own valuation basis.

\subsection{The haircut applies to the held value}

A lender's policy discount $h \in [0,1)$ applies to the held value,
$V_{\mathrm{lendable}} = V_{\mathrm{held}} (1 - h)$. This is not a matter of taste. A
haircut is calibrated to the volatility of a correctly measured value: it answers how far a
true value might fall before the lender can liquidate. Applying it to
$V_{\mathrm{reported}}$ instead gives
\[
  V_{\mathrm{reported}} (1 - h) = \big( V_{\mathrm{held}} + \Omega \big)(1 - h)
  = V_{\mathrm{lendable}} + \Omega (1 - h),
\]
so a residual $\Omega(1-h)$ of pure misstatement survives into the lender's book, scaled
down but never removed. Equivalently, the haircut that merely undoes the misstatement,
lending nothing against volatility at all, is
$h^{*} = 1 - V_{\mathrm{held}}/V_{\mathrm{reported}} = \delta/10^{4}$. A haircut absorbs
volatility. It does not absorb a misstatement.

Against $\mathrm{NAV}_{\mathrm{held}} = 0.803922$ and
$\mathrm{NAV}_{\mathrm{naive}} = 1.000000$, a pledge of $u = 1{,}000{,}000$ units:

\begin{center}
\begin{tabular}{lr}
\toprule
quantity & XRP \\
\midrule
$V_{\mathrm{reported}}$ & $1.000000$ \\
$V_{\mathrm{held}}$ & $0.803922$ \\
overstatement $\Omega$ & $\mathbf{0.196078}$ \\
\addlinespace
$V_{\mathrm{lendable}}$ at $h = 0.20$, correctly applied & $0.643138$ \\
haircut misapplied to $V_{\mathrm{reported}}$ & $0.800000$ \\
residual misstatement $\Omega(1-h)$ & $0.156862$ \\
\addlinespace
break-even haircut $h^{*}$ & $19.61\%$ \\
\bottomrule
\end{tabular}
\end{center}

A conventional $20\%$ haircut on the reported figure is, on this instrument, almost exactly
consumed by the misstatement alone. It leaves $0.39$ percentage points of protection against
everything else that could go wrong.

\section{Limitations}

\paragraph{Devnet only.} Every figure here was measured on the XRP Ledger Devnet, network id
$2$, against \texttt{rippled} 3.4.0-rc5. The lending amendments this depends on, XLS-65 for
vaults and XLS-66 for the broker and loan objects, are not active on Mainnet, so there is no
production facility to measure and no claim is made about one. The arithmetic is a property
of the protocol rules rather than of Devnet and should carry over unchanged if the
amendments activate. The hashes will not: Devnet is periodically reset.

\paragraph{A $29$ day history window.} The conduct score and the sequence score $\phi$ are
reconstructed from public transaction history, which the Devnet cluster retains for roughly
$29$ days. Behaviour before that window is invisible, so every conduct figure is a lower
bound on observable misconduct over a bounded window, not a statement about a manager's
whole record.

\paragraph{Par at $1.0$.} $\rho$ compares $A$ against $S$ and so assumes one unit was issued
per unit of asset subscribed. For a closed-ended vault this holds structurally rather than
by convention: subscription closes before any lending is permitted, so no profit or loss can
accrue between the first subscriber and the last, every subscriber enters at the same unit
value, and $S$ is exactly capital subscribed measured in asset units. \texttt{VaultKind} is
immutable, so a facility cannot leave that regime later. The assumption fails for an
open-ended vault, where subscribers enter at different unit values and $\rho$ would need the
full subscription cash-flow history rather than $S$ alone. No facility measured here can be
one: \texttt{LoanBrokerSet} returns \texttt{tecNO\_PERMISSION} against any vault that is not
closed-ended, so a vault that lends at all is in the regime by construction. An open-ended
vault would need that cash-flow history added before $\rho$ meant anything on it.

\paragraph{The closed form is not the implementation.} Lemma~\ref{lem:closed} drops the
ceiling and the per-step $\min$ under (A1) to (A3). Where a step is capped by the defaulting
principal or by remaining cover the closed form is no longer exact, although
Theorem~\ref{thm:order} stays directionally valid because a cap can only reduce a step. The
implementation therefore always evaluates the iterative recursion, which mirrors the ledger
step for step including the ceiling and both caps, and uses the closed form only as an
independent check. On the measured facility the two agree exactly, which is the point of
computing both.

\paragraph{Rounding, and the bands.} $\delta$ is an integer basis point rounded to nearest,
halves away from zero, so $1960.78$ is published as $1961$. One deviation is worth stating
outright: the project's internal specification writes that step with a floor, while the
implementation rounds, and every published figure follows the implementation. A floor would
report $1960$ where the ledger-facing code, the oracle object and the API all report $1961$.
Unit values are likewise rounded to nearest at the $6$ decimal places the valuation is
reported to, and oracle values at the dimension's scale, $6$ decimals for NAV and $4$ for the
other five. No figure here depends on a digit below those scales. The thresholds and
penalties in Section~5 are not derived from anything: they are a stated, ordinal,
reproducible judgement, published in full so a reader who disagrees can say where and by how
much. The two results that are not judgements are Proposition~\ref{prop:empty} and
Theorem~\ref{thm:order}, which follow from the protocol rules.

\end{document}
